\section{A globally convergent Euler--Hasse series}
\label{sec:hasse}

A further representation from the reference pages is available as an everywhere convergent series
of finite differences.  For $n\ge0$, define the entire finite sum
\[
 D_n(s)=\sum_{k=0}^n(-1)^k\binom nk(k+1)^{-s},\qquad
 \mathcal{E}_N(s)=\sum_{n=0}^N2^{-n-1}D_n(s).
\]
The point is to bound $D_n$ by a polynomial in $n$ on each bounded
complex region; the naive bound $2^n(n+1)^{|s|}$ would not suffice.

\subsection{A shifted finite Mellin identity}

Put $f_n(t)=e^{-t}(1-e^{-t})^n$ for $t\ge0$.  In the basis
\[
 g_{n,j}(t)=e^{-(j+1)t}(1-e^{-t})^{n-j},\qquad 0\le j\le n,
\]
finite differentiation gives
\[
 g_{n,j}'=-(j+1)g_{n,j}+(n-j)g_{n,j+1}.
 \tag{EH1}
\]
At $j=n$ the second term is zero and is omitted.  The sum of the
absolute values of the two displayed coefficients is $n+1$.
Since $f_n=g_{n,0}$, induction on $m$ proves that the absolute
coefficient sum of $f_n^{(m)}$ in this basis is at most $(n+1)^m$.
Each basis function has modulus at most $e^{-t}$; consequently
\[
 |f_n^{(m)}(t)|\le(n+1)^m e^{-t}\qquad(n,m\ge0).       \tag{EH2}
\]
This includes $n<m$; no exceptional initial terms are needed.

Let $P(z)=1/\Gamma(z)$ denote the entire reciprocal gamma construction.
For $\Re(s+m)>0$, positive scaling in Euler's integral gives
\begin{equation}
 \boxed{\displaystyle
 D_n(s)=(-1)^m P(s+m)
       \int_0^\infty t^{s+m-1}f_n^{(m)}(t)\,dt.}
 \label{eq:hasse-shifted-mellin}
\end{equation}
There is a direct finite proof, without analytic continuation: expand
\[
 f_n^{(m)}(t)=(-1)^m\sum_{k=0}^n
            (-1)^k\binom nk(k+1)^m e^{-(k+1)t}.
\]
Integrate this finite sum.  Its $k$th integral is
$\Gamma(s+m)(k+1)^{-s-m}$, and multiplication by $(k+1)^m$ and
$P(s+m)$ gives \eqref{eq:hasse-shifted-mellin}.
If instead one derives the identity by $m$ integrations by parts on
$\Re s>0$, the boundary terms are $t^{s+j}f_n^{(j)}(t)$.
At zero their moduli are at most $(n+1)^j t^{\Re s+j}$; at infinity
(EH2) bounds them by a polynomial times $e^{-t}$.  Thus all boundary
terms vanish with the gamma gateway's explicit cutoffs.

For actual evaluation of the integral, suppose
$\Re s\ge-A$, $|s|\le R$, where $A\ge0$ and $R\ge1$ are integers.
Choose $m=A+2$ and $S=R+m$.  At a lower endpoint $0<\epsilon\le1$
and an upper endpoint $T\ge1$, (EH2) gives
\[
 \int_0^\epsilon|t^{s+m-1}f_n^{(m)}(t)|\,dt
       \le(n+1)^m\epsilon^2/2,
\]
\[
 \int_T^\infty|t^{s+m-1}f_n^{(m)}(t)|\,dt
       \le(n+1)^m2^{S+1}S!e^{-T/2}.                 \tag{EH3}
\]
On the remaining finite interval an explicit Lipschitz bound is
\[
 (n+1)^mT^S\bigl((S+1)/\epsilon+n+1\bigr),
\]
by (EH1) and the product rule.  Apply the finite quadrature rule with
this bound.  These are endpoint and mesh certificates, not an appeal
to a qualitative improper integral.

Let $M_S$ be the explicit integer bound for $|P(z)|$ on $|z|\le S$
from the gamma gateway.  Since $2\le\Re s+m\le S$, the Euler
integral gives $\Gamma(\Re s+m)\le1+2^{S+1}S!$.  It follows that
\begin{equation}
 |D_n(s)|\le C_{A,R}(n+1)^m,
 \qquad C_{A,R}=M_S(1+2^{S+1}S!).
 \label{eq:hasse-polynomial-bound}
\end{equation}

\subsection{Global convergence and local power series}

The elementary factorial estimate
\[
 (n+1)^m\le m!8^m2^{n/2}
\]
holds for all $n,m\ge0$.  For $n\ge1$, use $n+1\le2n$,
$e^{n/4}\ge(n/4)^m/m!$, and $e\le4$; the case $n=0$ is immediate.
Geometric summation therefore proves
\begin{equation}
 \boxed{\displaystyle
 |\mathcal{E}(s)-\mathcal{E}_N(s)|\le 2C_{A,R}m!8^m\,2^{-N/2},
 \quad \mathcal{E}(s)=\sum_{n=0}^{\infty}2^{-n-1}D_n(s).}
 \label{eq:hasse-tail}
\end{equation}
Every inequality is first applied to a finite tail.  For example,
with $D=2C_{A,R}m!8^m$ and desired precision $2^{-p}$, take
$N=2(p+D+4)$.  The omitted tail is at most $2^{-p-4}$.
Evaluate each elementary number $(k+1)^{-s}$ to error at most
$2^{-p-4}/(N+1)$.  The binomial coefficient sum is $2^n$, so the
outer weight $2^{-n-1}$ makes the total evaluation error at most
$2^{-p-5}$.  Final rational rounding supplies the remaining error
reserve.  This is a fixed primitive recursive algorithm, although the
safe cutoffs are deliberately large.

For completeness, entireness also has a finite certificate.  Fix
$|s_0|\le R$ and use \eqref{eq:hasse-polynomial-bound} on
$|s|\le R+2$, taking $A=R+2$ there.  Write its constants as $C,m$.
The coefficient of $h^j$ in the finite entire function $D_n(s_0+h)$
is obtained by integrating its finite exponential Taylor polynomials
around $|h|=2$.  Fourier orthogonality and the factorial remainders give
coefficient modulus at most $C(n+1)^m2^{-j}$, by the finite circle
extraction already proved in the continuation argument.
Its weighted sum over $n$ therefore converges with the tail
\eqref{eq:hasse-tail} times $2^{-j}$.  The resulting $j$th Taylor
coefficient of $\mathcal{E}$ has modulus at most $2Cm!8^m2^{-j}$; on
$|h|\le1$ its degree-$J$ remainder is at most
$2Cm!8^m2^{-J}$.  Finite rectangles in $(n,j)$, with these two tail
bounds, prove equality with $\mathcal{E}(s_0+h)$.  Thus $\mathcal{E}$ is entire with
explicit Taylor data, rather than only a pointwise limit of entire
functions.

\subsection{Identification with eta and zeta}

For $\sigma=\Re s>0$, use the shifted integral with $m=0$ and sum
finitely over $n$.  The geometric identity is
\[
 \sum_{n=0}^N2^{-n-1}f_n(t)
 =\frac1{e^t+1}\left[1-
                  \left(\frac{1-e^{-t}}2\right)^{N+1}\right].
\tag{EH4}
\]
Therefore, if $I(s)=\int_0^\infty t^{s-1}/(e^t+1)\,dt$,
\[
 |\mathcal{E}_N(s)-P(s)I(s)|
       \le |P(s)|\Gamma(\sigma)2^{-N-1}.            \tag{EH5}
\]
Here $1/(e^t+1)\le e^{-t}$ gives the upper-endpoint estimate;
at zero one has the bound $\epsilon^\sigma/(2\sigma)$.
A supplied margin $\sigma\ge1/A$ makes both endpoints effective.
The finite alternating geometric identity also gives
\[
 \left|P(s)I(s)-\sum_{k=1}^K(-1)^{k-1}k^{-s}\right|
       \le |P(s)|\Gamma(\sigma)(K+1)^{-\sigma}.     \tag{EH6}
\]
Indeed the omitted kernel is $(-1)^K e^{-Kt}/(e^t+1)$, whose
absolute integral is at most
$\Gamma(\sigma)(K+1)^{-\sigma}$.  Both right sides have primitive
recursive moduli on the supplied positive half-plane chart.  Thus
$\mathcal{E}(s)=\eta(s)$ there.

To state the global identity without hiding a removable pole, put
$E(s)=(s-1)\zeta(s)$ with its entire Euler--Maclaurin definition, and
\[
 B(s)=\sum_{j=0}^{\infty}
       \frac{(-1)^j(\log2)^{j+1}(s-1)^j}{(j+1)!}.
\]
The factorial tail proves that $B$ is entire, and finite exponential
identities give $(s-1)B(s)=1-2^{1-s}$ and $B(1)=\log2$.
The entire functions $\mathcal{E}$ and $BE$ agree on $\Re s>1$, by the already
proved eta identity.  Apply the finite disk-propagation estimate to
their difference, using the explicit Taylor bounds above and those
for $B,E$.  It gives the entire identity
\begin{equation}
 \boxed{\displaystyle
 \mathcal{E}(s)=B(s)E(s),\qquad
 \eta(s)=(1-2^{1-s})\zeta(s),\qquad \eta(1)=\log2.}
 \label{eq:hasse-entire-eta}
\end{equation}
The middle expression means the entire continuation just defined,
including its removable value at $s=1$.

In particular, on every chart where $1-2^{1-s}$ has a supplied
positive rational modulus bound,
\[
 \boxed{\displaystyle
 \zeta(s)=\frac1{1-2^{1-s}}
 \sum_{n=0}^{\infty}\frac1{2^{n+1}}
       \sum_{k=0}^n(-1)^k\binom nk(k+1)^{-s}.}
\]
The numerator series converges everywhere.  The displayed quotient
is not used at its exceptional points
$s=1+2\pi i k/\log2$; at $k\ne0$ the entire numerator vanishes
and Euler--Maclaurin supplies the regularized zeta value.  This domain
qualification is essential for a constructive evaluation algorithm.

For completeness, the exceptional points also have a derivative formula.
For a nonzero integer $k$, put $s_k=1+2\pi i k/\log2$.
Differentiate the entire identity at $s_k$, where $s_k\mathrel\#1$.
The derivative of $1-2^{1-s}$ is $\log2$ there, so
\[
 \zeta(s_k)=\frac{\mathcal{E}'(s_k)}{\log2}.
\]
This is a licensed division by a positive constant and gives the same
regular value as the Euler--Maclaurin construction.
